\section{What is at stake? Three policy levers}\label{sec:stake}

The introduction set out two families of policy: if the invention is the problem, invest in research capability; if
the inventor is, relieve owners' constraints through credit, fee relief or markets for technology. The renewal model
can price the levers that the evidence points to on a common scale. I simulate each patent on its own fee path to
expiry with the deterministic model of Section~\ref{sec:values}, change one thing at a time, and report the change
in the returns owners earn from their patents per filing-year cohort, in 2024 PPP dollars and as a share of national
R\&D spending, for China and Korea\footnote{China and Korea are the two emerging-market
offices with fee discounts for small owners, which the first lever needs. Korea records no sales of granted patents,
so its markets-for-technology entries for firms use the premium estimated for China and Russia, and its universities
are not priced. Taiwan and the Czech Republic record no ownership transfers, the Czech Republic has no fee discount
for small owners, and the Russian values are not identified (Section~\ref{sec:values}).} (Table~\ref{tab:levers}).\footnote{\label{fn:levers}Patents a year: in China 4{,}400 of
small firms, 2{,}100 of mid-sized firms, 4{,}700 of large firms and 6{,}600 of universities; in Korea 5{,}600, 3{,}600,
26{,}200 and 3{,}300. Firm size: inventions filed worldwide in 2000--2005. National R\&D: gross domestic expenditure on
R\&D, 2000--2005 average \citep{WDI}. Decay rate $d = 0.85$; fees after 2024 at 2024 prices. Fee relief: every eligible
owner (firms with fewer than 100 inventions, universities) pays the discounted fees; the fees no longer paid, a
transfer not counted in the table, are \$8.6 million per cohort in China and \$1.9 million in Korea. Markets for technology:
returns multiplied, from one year after grant, by a factor calibrated to the group's extra post-sale survival in the
emerging-market offices (1.41 for small firms and 1.39 for mid-sized firms in China, 1.31 and 1.28 in Korea; 5.0 for
Chinese universities, calibrated to China's own university premium, since the advanced-economy one is slightly
negative); large firms show no extra post-sale survival. University patenting: mean log returns of Chinese
universities raised to those of foreign universities in China (upper bound) or to the foreign level plus the average
domestic--foreign university gap in the United States, Germany, the United Kingdom and France (lower bound).} Fees saved are transfers from the state to owners,
so they are not counted in the table.

\paragraph{Fee relief.} Suppose every eligible owner (firms with fewer than 100 inventions and universities) took up
the Chinese and Korean fee discounts. Owners of a Chinese cohort would pay \$8.6 million less in fees and owners of a
Korean one \$1.9 million less, but they would keep their patents only marginally longer, and the returns they earn would rise
by about half a million dollars in China and 1.3 million in Korea. Fees are small relative to what patents are worth,
so cutting them changes few renewal decisions.

\paragraph{Markets for technology.} Suppose that, from one year after grant, domestic owners' patents earned what they
would earn in the hands of a buyer. For each group I calibrate the buyer's advantage to the extra survival of patents sold by
emerging-market owners in Section~\ref{sec:constraints}: 6.1 points over five years for small firms and 5.0 for
mid-sized firms, relative to sellers of the same size in the advanced economies, and none for large firms, whose sold
patents do not gain more than advanced-economy firms' do. This requires returns about 40\% higher for small and
mid-sized firms in China and 30\% higher in Korea, where no sales of granted patents are recorded and the Chinese and
Russian premium is used. The returns from a cohort of firms' patents would rise by about \$390 million in China and
\$140 million in Korea, all of it from small and mid-sized firms. For Chinese universities, whose sold patents
survive 45 points longer, the calibrated advantage is a fivefold increase in returns and the gain \$710 million. Each figure is an upper bound, since it reallocates every patent of the group.

\paragraph{University patenting.} Suppose instead that Chinese universities' patents earned what foreign
universities' patents earn in China at equal citations. Each cohort would be worth about \$310 million more; with
the typical university home premium of the advanced-economy offices as the benchmark, about \$170 million. This and
the university row under markets for technology are two routes to the same lost value, improving what universities
patent or moving their patents to owners who use them, and should not be added.

\begin{table}[t]
\centering
\begin{threeparttable}
\caption{Returns at stake per filing-year cohort, by lever and owner group (millions of 2024 PPP dollars)}\label{tab:levers}
\small
% Source: em-patent-value output/policy_levers.csv (26_policy_levers.py)
\begin{tabular}{lcccccc}
\toprule
 & \multicolumn{2}{c}{Fee relief} & \multicolumn{2}{c}{Markets for technology} & \multicolumn{2}{c}{University patenting} \\
\cmidrule(lr){2-3}\cmidrule(lr){4-5}\cmidrule(lr){6-7}
 & China & Korea & China & Korea & China & Korea \\
\midrule
Small firms (1--9)          & 0.08 & 0.6 & 249 & 80 & -- & -- \\
Mid-sized firms (10--99)    & 0.03 & 0.4 & 146 & 63 & -- & -- \\
Large firms (100+)          & n.e. & n.e. & 0 & 0 & -- & -- \\
Universities                & 0.4  & 0.3 & 711 & n.a. & 174--310 & -- \\
\addlinespace
All owners                  & 0.5  & 1.3 & 1{,}106 & 143 & 174--310 & -- \\
\quad \% of national R\&D   & 0.00 & 0.00 & 1.17 & 0.35 & 0.19--0.33 & -- \\
\bottomrule
\end{tabular}
\begin{tablenotes}\footnotesize
\item \emph{Notes:} Change in the returns domestic owners earn from their patents, excluding fees, for one
filing-year cohort; deterministic renewal model, each patent simulated on its own fee path. Upper bounds for markets
for technology. The two university entries are alternative routes to the same value and are not additive. n.e.: not
eligible; n.a.: not available (Korea records no sales of granted patents). Details in footnote~\ref{fn:levers}.
\end{tablenotes}
\end{threeparttable}
\end{table}

\paragraph{Reading the numbers.} Three cautions apply. The figures are private returns to owners, not welfare: they
leave out knowledge spillovers and the gain to users when a patent enters the public domain. The deterministic model
gives conservative levels, several times below those of the learning model. And the markets-for-technology figures
are upper bounds. The gap between the levers survives all three: relieving owners of fees is worth about a million
dollars a cohort, while helping patents reach owners who can use them and improving what Chinese universities patent
are worth hundreds of millions, two to three orders of magnitude more.
